% !TEX root = ../1Main.tex
% Code-Beispiele im Anhang

\chapter{Code-Beispiele}
\label{chap:code-beispiele}

% ===================================================================
% ANLEITUNG: Dies ist ein Beispielanhang für Code-Listings.
% Diese Vorlage nutzt 'minted' für Syntax-Highlighting.
% WICHTIG: Kompilierung mit --shell-escape erforderlich!
% ===================================================================

Dieser Anhang demonstriert verschiedene Code-Listings mit dem modernen \texttt{minted}-Paket, das über 300 Programmiersprachen unterstützt.

\section{Python-Beispiel}
\label{sec:python-code}

Listing \ref{lst:python-beispiel} zeigt ein einfaches Python-Programm mit Syntax-Highlighting.

\begin{listing}[htbp]
\begin{minted}[
    linenos,
    frame=lines,
    framesep=2mm,
    baselinestretch=1.2,
    bgcolor=gray!10,
    fontsize=\footnotesize
]{python}
def fibonacci(n):
    """
    Berechnet die n-te Fibonacci-Zahl rekursiv.
    
    Args:
        n: Index der Fibonacci-Zahl
        
    Returns:
        Die n-te Fibonacci-Zahl
    """
    if n <= 1:
        return n
    return fibonacci(n-1) + fibonacci(n-2)

# Hauptprogramm
if __name__ == "__main__":
    for i in range(10):
        print(f"F({i}) = {fibonacci(i)}")
\end{minted}
\caption{Python-Implementierung der Fibonacci-Funktion}
\label{lst:python-beispiel}
\end{listing}

\section{Java-Beispiel}
\label{sec:java-code}

Listing \ref{lst:java-beispiel} demonstriert objektorientierte Programmierung in Java.

\begin{listing}[htbp]
\begin{minted}[
    linenos,
    frame=lines,
    framesep=2mm,
    baselinestretch=1.2,
    bgcolor=gray!10,
    fontsize=\footnotesize
]{java}
public class HelloWorld {
    private String message;
    
    /**
     * Konstruktor
     */
    public HelloWorld(String message) {
        this.message = message;
    }
    
    /**
     * Gibt die Nachricht aus
     */
    public void printMessage() {
        System.out.println(message);
    }
    
    public static void main(String[] args) {
        HelloWorld hello = new HelloWorld("Hello, World!");
        hello.printMessage();
    }
}
\end{minted}
\caption{Java Hello World mit Klassen}
\label{lst:java-beispiel}
\end{listing}

\section{Scala-Beispiel}
\label{sec:scala-code}

Listing \ref{lst:scala-beispiel} zeigt funktionale Programmierung in Scala.

\begin{listing}[htbp]
\begin{minted}[
    linenos,
    frame=lines,
    framesep=2mm,
    baselinestretch=1.2,
    bgcolor=gray!10,
    fontsize=\footnotesize
]{scala}
// Funktionale Programmierung in Scala
object FunctionalExample {
  // Higher-order functions
  def applyTwice(f: Int => Int, x: Int): Int = {
    f(f(x))
  }
  
  // Pattern Matching
  def describe(x: Any): String = x match {
    case i: Int if i > 0 => "positive integer"
    case i: Int if i < 0 => "negative integer"
    case 0 => "zero"
    case s: String => s"string: $s"
    case _ => "something else"
  }
  
  def main(args: Array[String]): Unit = {
    val increment = (x: Int) => x + 1
    println(applyTwice(increment, 5))  // Output: 7
    println(describe(42))              // Output: positive integer
  }
}
\end{minted}
\caption{Scala Beispiel mit funktionaler Programmierung}
\label{lst:scala-beispiel}
\end{listing}

\section{JSON-Datei}
\label{sec:json-code}

Auch Konfigurationsdateien können formatiert werden:

\begin{listing}[htbp]
\begin{minted}[
    linenos,
    frame=lines,
    framesep=2mm,
    baselinestretch=1.2,
    bgcolor=gray!10,
    fontsize=\footnotesize
]{json}
{
  "name": "example-project",
  "version": "1.0.0",
  "dependencies": {
    "react": "^18.0.0",
    "typescript": "^5.0.0"
  },
  "scripts": {
    "build": "tsc",
    "test": "jest"
  }
}
\end{minted}
\caption{package.json Beispiel}
\label{lst:json-beispiel}
\end{listing}

\section{Shell-Script}
\label{sec:bash-code}

\begin{listing}[htbp]
\begin{minted}[
    linenos,
    frame=lines,
    framesep=2mm,
    baselinestretch=1.2,
    bgcolor=gray!10,
    fontsize=\footnotesize
]{bash}
#!/bin/bash
# Kompilierungs-Script für LaTeX

echo "Starting LaTeX compilation..."

# Erste Kompilierung
lualatex --shell-escape 1Main.tex

# Glossaries erstellen
makeglossaries 1Main

# Bibliographie erstellen
biber 1Main

# Erneute Kompilierungen für Referenzen
lualatex --shell-escape 1Main.tex
lualatex --shell-escape 1Main.tex

echo "Compilation complete!"
\end{minted}
\caption{Bash-Script für LaTeX-Kompilierung}
\label{lst:bash-beispiel}
\end{listing}

% ===================================================================
% HINWEIS: Fügen Sie Ihren eigenen Code hier ein!
% Unterstützte Sprachen: Python, Java, C, C++, JavaScript, Scala,
% Rust, Go, Ruby, PHP, SQL, HTML, CSS, und viele mehr!
% ===================================================================